% The step estimator, the at-cost instrument Zb, the selection control and the null band.
% Build: latexmk -pdf at-cost-instrument.tex   (from docs/derivations/)
\documentclass[11pt]{article}
\usepackage[utf8]{inputenc}
\usepackage[T1]{fontenc}
\usepackage[a4paper,margin=2.5cm]{geometry}
\usepackage{amsmath,amssymb,amsthm}
\usepackage{booktabs,array}
\usepackage{textcomp}
\usepackage[hidelinks]{hyperref}
\usepackage{microtype}

\newtheorem*{assumption}{Assumption}
\newcommand{\E}{\mathbb{E}}
\newcommand{\Xt}{\tilde X}
\newcommand{\eur}{\text{\texteuro}}
\newcommand{\code}[1]{\texttt{\small #1}}

\title{The step estimator and the at-cost instrument}
\author{Spanish electricity market project}
\date{27 October 2026}

\begin{document}
\maketitle

\noindent This note writes out what \code{step\_engine.py}, \code{leg\_b\_levels.py} and
\code{price\_taking\_simulation.py} compute, and maps each symbol to the data. The scripts
are the authority: where this note and the code disagree, the code is right.

\section{A worked example}

A firm offers a hydro step at \eur62 per MWh, and its water value is \eur55, so the markup
is \eur7. In 4 of the 40 simulated hours this step sets the price:

\begin{center}
\begin{tabular}{ccc}
\toprule
draw & firm's net position $n$ (MW) & rivals' slope $s$ (MW per \eur) \\
\midrule
1 & 300 & 100 \\
2 & 200 & 50 \\
3 & 100 & 50 \\
4 & 200 & 200 \\
\bottomrule
\end{tabular}
\end{center}

\noindent The pooled incentive is $\Xt = (300+200+100+200)/(100+50+50+200) = 800/400 =
\eur2$. A price-taker would offer at \eur55 (markup 0); a firm exercising full market power
would offer at \eur57 (markup $\Xt$). The parameter $\beta$ is the share of $\Xt$ that
appears in markups, across many steps.

\section{The model}

\paragraph{Setting.} Firm $f$ submits a step curve for hour $H$ before it knows its residual
demand. That uncertainty is represented by $R=40$ equally likely draws $r=1,\dots,R$ of the
rivals' curve (Section~\ref{sec:data}). Let $RD_r(p)$ be $f$'s residual demand in draw $r$:
market demand minus rivals' supply at price $p$. At the clearing price $p_r^*$, $f$'s net
position (sales minus purchases, exposed volume only) is
\[
  n_r = RD_r(p_r^*),
\]
and the slope of its residual demand is
\[
  s_r = \lvert RD_r'(p_r^*)\rvert \approx
  \frac{\text{rivals' MW offered within } \pm h \text{ of } p_r^*}{2h}, \qquad h = \eur5.
\]

\paragraph{The draws.} Take $H$ = 20:00 on Tuesday 14 March 2023. The firm bids at noon on
13 March without knowing how rivals will bid. The candidates are the 20:00 hour of each day
from 14 February to 11 April, excluding 14 March itself: 56 evenings. Of these, the 40 whose
forecast of net load (REE's D$-1$ demand forecast minus its wind and solar forecasts) is
nearest 14 March's are kept. In each kept evening the market's real curves are taken, $f$'s
own offers removed, and the rest shifted sideways by the difference in forecasts:
\[
  \Delta = (1-d_f)\,\Delta F_{\text{demand}} - (1-w_f)\,\Delta F_{\text{wind}} - (1-s_f)\,\Delta F_{\text{solar}},
\]
with $d_f$, $w_f$, $s_f$ the firm's shares of buy offers, wind and solar. $f$'s actual curve for
14 March is then cleared against each of the 40 rival curves. The idea is that how rivals
bid on similar evenings stands in for $f$'s uncertainty about how they will bid on this one.
Its weakness: the spread across evenings mixes true uncertainty with real changes between
days (a move in the gas price from 9 to 14 March), so it may overstate that uncertainty.

\paragraph{One step.} Step $k$ of unit $u$ is offered at price $q_k$; the unit's marginal
(opportunity) cost is $c_k$. Let
\[
  M_k = \{\, r : p_r^* = q_k \,\}
\]
be the draws in which step $k$ sets the price. In the other draws a small change in $q_k$
does not move the price, so it does not move profit.

\paragraph{First-order condition.} In a draw $r \in M_k$, profit as a function of the price
is $\pi_r(p) = (p - c_k)\,RD_r(p)$ near $p=q_k$, plus terms that do not depend on $p$. So
\[
  \pi_r'(q_k) = RD_r(q_k) + (q_k - c_k)\,RD_r'(q_k) = n_r - (q_k - c_k)\,s_r .
\]
Raising the price by one euro earns one euro on each of the $n_r$ net MW and loses the
margin $q_k - c_k$ on the $s_r$ MW that rivals take. The step's price is optimal in
expectation when
\[
  \sum_{r \in M_k} \bigl[\, n_r - (q_k - c_k)\, s_r \,\bigr] = 0
  \quad\Longleftrightarrow\quad
  q_k - c_k = \frac{\sum_{r \in M_k} n_r}{\sum_{r \in M_k} s_r} \equiv \Xt_k .
\]
This is the condition of Wolak (2007), Kastl (2011) and Reguant (2014). $\Xt_k$ is a ratio of
sums, not a mean of the ratios $n_r/s_r$.

\paragraph{Conduct.} Nesting price-taking ($\beta=0$) and full market power ($\beta=1$):
\begin{equation}\label{eq:est}
  y_k \equiv q_k - c_k = \beta\,\Xt_k + \alpha_{u,m} + \varepsilon_k ,
\end{equation}
where $\alpha_{u,m}$ is a unit $\times$ month fixed effect (unit $\times$ day in some rows) and
$\varepsilon_k$ is everything unobserved, including any error in $c_k$. Observations are steps
that set the price in at least one draw ($\lvert M_k\rvert > 0$), weighted by $\lvert M_k
\rvert$, with a mean slope $\sum_{M_k} s_r / \lvert M_k\rvert$ of at least 20 MW per \eur.
For hydro, $c_k = \lambda^*_{30s}$, the water value, and $y_k$ is \code{gap30}; with unit
$\times$ day effects the daily cost is absorbed and $y_k$ is the offer price.

\section{Why \texorpdfstring{$\Xt$}{X} needs an instrument}

$\Xt_k$ is built on the firm's own curve. Both $M_k$ (in which draws step $k$ sets the price)
and $n_r$ (the firm's position there) depend on the firm's own offer prices, and those
contain $\varepsilon$. A step priced high because its true cost is high sets the price in
different, higher-priced draws than a step priced low. So $\operatorname{Cov}(\Xt_k,
\varepsilon_k) \neq 0$, and least squares on \eqref{eq:est} is biased.

\section{The at-cost instrument \texorpdfstring{$Z^b$}{Zb}}

\paragraph{Counterfactual curve.} Move every one of $f$'s storage and pumped-storage hydro
steps priced above \eur1 to its unit's water value, $q_j \to c_j$, and leave every other
offer as submitted, including $f$'s own non-hydro offers. Re-clear each draw against the
same rivals' curves, giving $\tilde p_r$, $\tilde n_r$ and $\tilde s_r$.

\paragraph{Cost point and near draws.} Let $\bar c_d$ be the median water value across hydro
units on day $d$, and $N_k$ the 5 draws whose $\tilde p_r$ is nearest $\bar c_d$. Then
\[
  Z^b_k = \frac{\sum_{r \in N_k} \tilde n_r}{\sum_{r \in N_k} \tilde s_r},
\]
defined when $\sum_{N_k} \tilde s_r / 5$ is at least 20 MW per \eur. It is $\Xt$ evaluated
where the firm is a price-taker in its hydro, that is, at $\varepsilon = 0$: the
approximation to optimal instruments of Berry, Levinsohn and Pakes (1999). Using the class
median $\bar c_d$ rather than the unit's own $c_k$ keeps an error in one unit's cost out of
the instrument; $Z^c$ uses $c_k$ itself.

\begin{assumption}[Exclusion]
$\E[\varepsilon_k \mid Z^b_k, \alpha_{u,m}] = 0$.
\end{assumption}

\noindent $Z^b$ contains neither $q_k$ nor $c_k$. It moves with rivals' curves, the forecasts
that select and shift the draws, $f$'s non-hydro offers and $\bar c_d$. Relevance:
$\operatorname{corr}(\Xt, Z^b) = 0.44$, first-stage $F$ about 90 to 410 across specifications.

\paragraph{What would violate it.}
\begin{enumerate}
  \item The unobserved cost moves with the shape of rivals' curves: a firm values its water
    more on days when rivals' supply near the price is thin. The planted test
    (\code{--sim-scale}: price-takers at $m_f\lambda^*$, with $m_f$ each firm's ratio of
    offer to $\lambda^*$) is one check; it moves $\beta$ by $-0.01$ to $-0.05$.
  \item $f$'s own non-hydro offers are kept as submitted, so $Z^b$ is not wholly free of
    $f$'s choices.
\end{enumerate}

\paragraph{A wrong cost level does not bias $\beta$.} If $c_k = c_k^{\text{true}} +
\delta_{u,m}$, the error is constant within a unit-month and is absorbed by $\alpha_{u,m}$. For
the same reason the test cannot confirm the cost level. That rests on separate checks:
EDP's revealed water value is within 6\% of $\lambda^*$.

\section{The selection control}

A step enters the sample only if some draw's price reaches it, so the kept steps' markups
rise with how far the hour's prices sit above the unit's cost, and the instruments move with
that too. Let $\bar{\tilde p}_k$ and $\sigma_k$ be the mean and standard deviation of the
draws' prices on the at-cost curve, and
\[
  P^{\text{rel}}_k = \bar{\tilde p}_k - c_k .
\]
The corrected specification gives each unit $i$ its own slope:
\[
  y_k = \beta\,\Xt_k + \gamma_i\,P^{\text{rel}}_k + \alpha_{u,m} + \varepsilon_k ,
\]
with $\Xt_k$ instrumented by $Z^b_k$. The variants add $\bigl(P^{\text{rel}}_k\bigr)^2/100$ or
$\sigma_k$, each also with a slope per unit; \code{none} drops the control. The draws' prices
are taken on the at-cost curve, so the firm's own offers do not enter the control either.

\section{The null band}

\paragraph{Simulated price-takers.} For seed $s$, replace $f$'s hydro steps by $c_j +
\eta_j^{(s)}$, where $\eta^{(s)}$ is the unit's own markups reshuffled, which ignores $f$'s
stake in the price. Then $\beta = 0$ by construction. Rebuild $\Xt$, $Z^b$ and the sample on
the same draws and estimate $\hat\beta^{(s)}$.

\paragraph{Reading rule} (fixed 2026-10-08, before the actual offers were read). With the
\code{level} control,
\[
  B = \Bigl[\, \min_{s \ge 1} \hat\beta^{(s)},\ \max_{s \ge 1} \hat\beta^{(s)} \,\Bigr].
\]
The actual offers ($s=0$) differ from price-takers only if their Anderson--Rubin 95\% set
does not meet $B$. Full market power is rejected if 1 lies outside that set. More seeds can
only widen $B$, so a ``differs'' on two seeds is provisional; seeds 3 to 10 are the next run.

\section{Where each symbol comes from}\label{sec:data}

\begin{center}
\small
\begin{tabular}{>{\raggedright}p{2.6cm} >{\raggedright}p{6.4cm} >{\raggedright\arraybackslash}p{5.2cm}}
\toprule
symbol & data & file or code \\
\midrule
$q_k$, $f$'s curve & OMIE day-ahead bid curves (\code{curva\_pbc\_uof}); the effective curve: accepted rows below $p^*$, offered rows at and above it & \code{step\_curves.py} $\to$ \code{own\_}, \code{market\_YYYYMM} \\
$f$: whose units & OMIE unit listings with ownership shares, dated & \code{lib/unit\_ownership.py}, \code{owner\_groups.parquet} \\
draws $r$ & 40 hours at the same clock hour within $\pm 28$ days, nearest in REE's D$-1$ forecast of net load (ESIOS 1775, 1777, 1779), shifted to $H$'s forecasts & \code{step\_engine.py} \\
$n_r$ & $f$'s exposed sales minus purchases at $p_r^*$; regulated pass-through suppliers excluded & \code{lib/omie\_owner.py} \\
$s_r$ & rivals' MW within $\pm$\eur5 of $p_r^*$, over \eur10; undefined below 5 MW per \eur & \code{step\_engine.py} \\
$c_k$, hydro & $\lambda^*_{30s}$: OMIP forwards $\times$ last year's hourly shape, the unit's running share, MITECO or SNIRH reservoir fill; known at D$-2$ & \code{hydro\_cost\_daily.parquet} \\
$c_k$, thermal & MIBGAS gas and EUA carbon (coal: API2 over the ECB rate) divided by JRC efficiency; prices final at noon on D$-1$ & \code{implied\_marginal\_\allowbreak cost\_daily.parquet}, \code{\_bid} columns \\
$\sum_{M_k} n_r$, $\sum_{M_k} s_r$, $\lvert M_k\rvert$ & & \code{hmoments\_YYYYMM}: \code{m\_sn}, \code{m\_s}, \code{m\_n} \\
$\sum_{N_k} \tilde n_r$, $\sum_{N_k} \tilde s_r$ & & \code{hcompcost\_YYYYMM}: \code{kb\_sn}, \code{kb\_s} \\
$\bar{\tilde p}_k$, $\sigma_k$ & & \code{hcompcostprice\_YYYYMM}: \code{cp\_mean}, \code{cp\_sd} \\
$\hat\beta^{(s)}$ inputs & simulated price-takers, seeds 1--3 & \code{hsim}$s$\code{\_moments}, \code{hsim}$s$\code{\_units} \\
planted test & price-takers at $m_f \lambda^*$ & \code{hsimw1\_*} \\
\bottomrule
\end{tabular}
\end{center}

\noindent Thermal uses the same caches without the \code{h} prefix, with $f$'s gas and coal
steps in place of its hydro steps.

\paragraph{What the October rebuild changes.} Only the membership of $f$. ADOURO's Douro dams
are Movhera's, not EDP's, from 2020-12-17 to 2021-09-08, and Castejón's represented units are
TotalEnergies', not EDP's, from February 2024 to January 2025. In those months the own curves
of EDP, Movhera and TotalEnergies change, and with them $n_r$, $M_k$, the rivals' curves,
$\Xt$, $Z^b$ and $P^{\text{rel}}$. Every other firm's rows were checked to be identical. The
costs and the formulas above are unchanged.

\end{document}
